% Standalone wrapper to render federated_method.tex as a self-contained PDF.
% Compile with:  tectonic federated_method_standalone.tex
\documentclass[11pt]{article}

\usepackage[T1]{fontenc}
\usepackage[utf8]{inputenc}
\usepackage[margin=1in]{geometry}
\usepackage{amsmath}
\usepackage{amsfonts}
\usepackage{amssymb}
\usepackage{booktabs}
\usepackage{multirow}
\usepackage{array}
\usepackage{ragged2e}
\usepackage{amsthm}
\usepackage{algorithm}
\usepackage{algpseudocode}
\usepackage{lmodern}     % provides bold small-caps shape (silences T1/lmr/bx/sc warning)
\usepackage{microtype}  % character protrusion: better justification, fewer bad boxes

\emergencystretch=3em  % absorb minor overfull boxes in this draft preview

\theoremstyle{plain}
\newtheorem{proposition}{Proposition}
\theoremstyle{definition}
\newtheorem{remark}{Remark}

% Marks a pre-registered result that has not yet been measured (see convergence protocol).
\newcommand{\pending}{\textcolor{gray}{\emph{[to run]}}}
\usepackage{xcolor}

\title{Federating a Discrete Latent Vocabulary for\\Univariate KPI Anomaly Detection\\
\large Sufficient-Statistic Codebook Merging: a Benchmark and Mechanism Analysis}
\author{}
\date{}

\begin{document}
\maketitle

\begin{abstract}
We study how to federate TimeVQVAE-AD across \emph{univariate} clients---one KPI per
client, which never share raw signals---for time-series anomaly detection. The detector
learns a discrete time--frequency representation with a VQ-VAE (Stage~1) and fits a
MaskGIT masked prior over the resulting code indices (Stage~2). The obstacle is that a VQ
codebook is not a gradient-trained parameter and its entries carry no fixed cross-client
meaning, so naive weight-averaging is the wrong aggregation primitive. We study a two-part
recipe: \textbf{(A)}~merge the global codebook from per-code \emph{sufficient
statistics}---assignment counts and centroid-sums---reproducing the pooled $k$-means
M-step (federated $k$-means; we claim no new theorem, only that this is the correct
primitive for a VQ codebook, vs.\ the weight-\textsc{FedAvg} of prior federated-codebook
work); and \textbf{(B)}~partially personalize the prior, \textsc{FedAvg}-ing its
transformer body while keeping small per-client output heads local.

Our benchmark is real: WSD, $31$ univariate web-service KPIs (AnoTransfer), grouped into
$4$ train-only cohorts; a controlled univariate synthetic benchmark with an explicit
heterogeneity axis serves for mechanism analysis. We report an \emph{honest, in-progress}
% The 0.31-vs-0.33 calibration claim was WITHDRAWN 2026-07-16: 0.33 has no artifact and 0.31
% was the under-trained pilot. See federated_method.tex:253. Baseline is pre-registered, pending.
picture. \textbf{(1)~Calibration:} a mandatory moving-average calibration baseline is
pre-registered and not yet measured, so we make no claim about the detector's absolute
advantage; all claims are relative (arm vs arm under one protocol) by construction.
\textbf{(2)~Pilot viability:} at a short budget the federated recipe is
the worst of local / centralized / federated in every cohort ($\Delta_{\text{fed}-\text{local}}
={-}0.19$ VUS-PR, $n{=}31$); but every arm is far from convergence, and on univariate
synthetic data the sufficient-statistic-vs-\textsc{FedAvg} contrast is exactly zero at this
budget, so this pilot is \emph{not} evidence about the method. We pre-register the
converged confirmatory sweep. \textbf{(3)~Mechanism:} because the deployed KPI score is the
prior term alone, any federated collapse is a \emph{prior} collapse (the Stage-1 codebook
stays healthy); we localize it and make it CPU-verifiable from persisted checkpoints.
\textbf{(4)~When federation hurts:} on synthetic data the federated--local gap is monotone
in cohort heterogeneity, but on real data that signal is null---the only predictor is a
client's own detectability, i.e.\ federation acts as regression-to-the-mean.

We are explicit about what is and is not established: the sufficient-statistic primitive's
advantage, and whether federation reaches local/centralized parity, are \emph{open}
questions resolved only by the pre-registered converged experiment; the calibration,
mechanism, and heterogeneity findings above are the current, real-data contributions.
\end{abstract}

\section{Introduction}
\label{sec:fed-intro}

Time-series anomaly detection underpins monitoring in web services, telemetry,
industrial plants, and cyber-physical systems. In many of these settings the data are
partitioned across \emph{clients}---web-service KPIs, operators, hospitals, utility
substations---each a \emph{single univariate} signal (one key performance indicator)
instrumented under different operating conditions.
Pooling these series into a single training corpus is often impossible: raw signals
are commercially sensitive or regulated, and moving them off-premises is precisely
what data-governance policies forbid. \emph{Cross-client federated learning} (FL)
answers this by training a shared model while raw data never leave a
client~\cite{mcmahan2017fedavg,kairouz2021advances}. It is a natural fit for anomaly
detection, where each client has abundant local \emph{normal} behavior but cannot
export the signals themselves.

We build on TimeVQVAE-AD, a detector that (Stage~1) learns a
discrete time--frequency representation by quantizing an STFT with a
VQ-VAE~\cite{van_den_oord2017neural,lee2023timevqvae,lee2024timevqvaead} and
(Stage~2) fits a MaskGIT masked prior over the resulting code
indices~\cite{chang2022maskgit}; detection combines the prior's masked-NLL surprise
with a reconstruction-energy term. This discrete, two-stage design is attractive, but
federating it is an open problem. Discrete/codebook representations \emph{have} been federated
before---for classification, generative synthesis, and semantic
communication~\cite{zhang2024uefl,jiang2026synthfed,park2022fcdeepsc,jin2025svqvae}---and
federated time-series anomaly detection is itself emerging~\cite{perales2023temporalfed};
but to our knowledge no established recipe federates a VQ-based anomaly
\emph{detector}, in which the code \emph{indices} are themselves consumed by a
downstream masked prior, and where existing federated codebooks are either
personalized per client or averaged as weights for a reconstruction/transmission
objective (\S\ref{sec:related-work}). The obstacle is that generic
FedAvg~\cite{mcmahan2017fedavg} assumes gradient-trained, semantically aligned
parameters that can be averaged elementwise, and a VQ codebook violates both
assumptions. It is overwritten each step from EMA
assignment statistics rather than by a loss gradient, and its entries form an
\emph{unordered} set whose index $j$ has no fixed meaning across clients. Averaging
codebook weights across clients is therefore not the pooled update it appears to be,
and---as we show---it silently corrupts the very signal the detector depends on.

This paper proposes and validates a two-part federation recipe for this
architecture, and gives an honest, mechanism-level account of what works and what
does not. \textbf{Contribution~A} (\S\ref{sec:fed-stage1}) federates the Stage-1
codebook by aggregating per-code \emph{sufficient statistics}---assignment counts
$n_j$ and centroid-sums $m_j$---rather than weights. Given a frozen broadcast
codebook, the server's merge \emph{exactly reproduces the pooled $k$-means/EM
M-step for the round's assignments} (Proposition~1, with a unit test); it is the
aggregation \emph{operator} that is exact, not an equivalence of the federated and
centralized training processes. The pooled M-step is itself the standard distributed
Lloyd step of Dhillon \& Modha~\cite{dhillon1999parallel}---same $(n_j,m_j)$ accumulators,
same sum-allreduces, same smoothed denominator---equivalently the VQ-VAE EMA
update~\cite[App.~A.1]{van_den_oord2017neural} with the accumulation boundary moved from
the minibatch to the client; the one-shot k-FED~\cite{dennis2021kfed} is a different
algorithm and not this precedent. We therefore claim the operator neither as new nor as
``the correct primitive for a VQ codebook'': Dhillon \& Modha's descent guarantee assumes
every processor applies the same distance map in a fixed $\mathbb{R}^d$, whereas our E-step
is a per-client \emph{learned} encoder that drifts each round, so the exactness is real but
descends no single objective. Our contribution is the measurable consequence---the merged
codebook is exact to fp32 while the client token languages it induces have disjoint
support---and the empirical finding that weight-averaging breaks not reconstruction but the
\emph{prior}. \textbf{Contribution~B} (\S\ref{sec:fed-stage2}) makes the Stage-2
prior \emph{partially personalized}: the transformer body is FedAvg'd, while the
channel positional embedding and per-position output bias stay local. This is a
domain-grounded instance of partial model
personalization~\cite{arivazhagan2019fedper,collins2021fedrep,pillutla2022partial,li2021fedbn},
identifying which prior components encode client-specific normality.

We are deliberately honest about what is and is not yet established (full numbers in
\S\ref{sec:fed-results}). On the \emph{real} univariate KPI benchmark we report three
current, real-data findings and one pre-registered open question. \emph{Calibration:}
% WITHDRAWN 2026-07-16 (0.33 unbacked, 0.31 = pilot); see federated_method.tex:253.
a mandatory moving-average calibration baseline is pre-registered and \pending, so every claim
in this paper is \emph{relative} (arm vs arm under one protocol), never absolute superiority.
\emph{Mechanism:} because the deployed KPI score is the prior's masked-NLL term alone,
any federated collapse is by construction a \emph{prior} collapse---the Stage-1 codebook
stays healthy---consistent with the hypothesis that weight-\textsc{FedAvg} of the codebook
breaks the prior by de-aligning code \emph{indices} across clients. \emph{When federation
hurts:} on controlled synthetic data the federated--local gap is monotone in cohort
heterogeneity, but on real data that signal vanishes and the sole predictor is a client's
own detectability, i.e.\ federation behaves as regression-to-the-mean. The \emph{open}
question, deferred to a pre-registered converged experiment, is whether the
sufficient-statistic merge beats weight-\textsc{FedAvg}---and whether federation reaches
local/centralized parity---at convergence: a short-budget pilot leaves every arm
under-trained and is not evidence about the method. Our reporting follows ML-for-science
standards~\cite{kapoor2024reforms}, and this study complements recent federated
time-series anomaly detection~\cite{perales2023temporalfed}.

\paragraph{Contributions.}
\begin{itemize}
  \item \textbf{(A) Codebook by sufficient statistics (\S\ref{sec:fed-stage1}).} A
  federation rule for the VQ codebook that aggregates per-code counts and
  centroid-sums; given the frozen broadcast codebook, its server-side merge
  reproduces the exact pooled $k$-means/EM M-step for the round (Proposition~1),
  keeping code indices globally coherent across clients.
  \item \textbf{(B) Partially-personalized prior (\S\ref{sec:fed-stage2}).} A
  domain-grounded instance of partial model personalization for the MaskGIT prior:
  the transformer body is shared and FedAvg'd, while the channel positional
  embedding and per-position output bias remain local.
  \item \textbf{A real-data benchmark and an honest, in-progress evaluation
  (\S\ref{sec:fed-results}).} On $31$ real univariate KPIs we establish:
  % withdrawn: "a mandatory moving-average calibration (the deep model does not beat it)" --
  % that baseline is unmeasured, see federated_method.tex:253.
  a mechanism localization (the deployed score is the prior alone, so federated collapse
  is a prior collapse with a healthy codebook, CPU-verifiable from checkpoints); and a
  heterogeneity analysis (the synthetic ``heterogeneity drives the gap'' law does not hold
  on real data---detectability does). The primitive's advantage and local/centralized
  parity are pre-registered as the converged confirmatory experiment, not claimed here.
\end{itemize}

\section{Related Work}
\label{sec:related-work}

\paragraph{Discrete-representation generative time-series modeling and anomaly detection.}
Vector-quantized autoencoders~\cite{van_den_oord2017neural} learn a discrete latent
codebook, turning a continuous signal into a sequence of code indices that a
downstream model can treat as tokens. TimeVQVAE~\cite{lee2023timevqvae} ports this
idea to time series by quantizing a time--frequency (STFT) representation, and
TimeVQVAE-AD~\cite{lee2024timevqvaead} builds an anomaly detector on top: a MaskGIT
masked prior~\cite{chang2022maskgit} is fit over the discrete tokens, and detection
combines the prior's masked negative-log-likelihood (token-level surprise) with a
reconstruction-energy term. Discrete/VQ representations have been used for anomaly
detection more broadly---e.g.\ hierarchical VQ transformers for
images~\cite{lu2023hvqtrans} and quantized autoencoders for multivariate 5G
telemetry~\cite{trappolini2025qaed}---but, like TimeVQVAE-AD, always in the
\emph{centralized} setting. We adopt this two-stage discrete pipeline unchanged at
the single-client level. Our concern is orthogonal to its modeling choices: how to
learn the shared codebook (Stage~1) and the masked prior (Stage~2) when the
training data are partitioned across clients that cannot exchange raw signals.

\paragraph{Federated and personalized learning.}
FedAvg~\cite{mcmahan2017fedavg} aggregates client models by weight averaging and is
the natural first attempt for both of our stages---yet it is the wrong primitive for
a VQ codebook. Our Stage-1 merge from per-code sufficient statistics (counts and
centroid-sums; Proposition~1, \S\ref{sec:fed-stage1}) is not new as an
\emph{algorithm}: it is the classical one-shot federated $k$-means / EM
M-step~\cite{dennis2021kfed}, and
it mirrors the prototype-aggregation view of FedProto~\cite{tan2022fedproto} applied
to VQ codes. What is new is its \emph{application} to a VQ codebook, together with
the empirical hypothesis (\S\ref{sec:fed-results}) that
explains why it should matter here: weight-FedAvg of the codebook does not harm
reconstruction---the Stage-1 codebook stays healthy---but it
silently misaligns code \emph{indices} across clients, and because the prior consumes
indices, this collapses the downstream prior anomaly signal. The sufficient-statistic
merge keeps index assignment globally coherent because it reproduces the pooled
M-step for the round.

For Stage~2 we keep a shared transformer body (FedAvg'd) while retaining the channel
positional embedding and per-position output bias \emph{locally}. This is partial
model personalization, a well-studied family: client-specific heads or
layers~\cite{arivazhagan2019fedper,collins2021fedrep}, the general
partial-personalization framework of~\cite{pillutla2022partial}, local
batch-normalization statistics~\cite{li2021fedbn}, and regularized or
Moreau-envelope personalization~\cite{li2021ditto,dinh2020pfedme}. We propose no new
personalization mechanism; our contribution (B, \S\ref{sec:fed-stage2}) is the
domain-grounded identification of \emph{which} prior components encode client-specific
normality (channel embedding and output bias), so that the shared masked prior can
be globally consistent while still respecting per-client period/jitter heterogeneity.

\paragraph{Federated discrete and codebook representations.}
A distinct line \emph{does} federate discrete/codebook representations, but outside
anomaly detection and with different design goals, so none supplies a recipe for the
problem we target. Closest to our Stage~1, UEFL~\cite{zhang2024uefl} federates a
discrete codebook across heterogeneous data clients---yet for classification, and by
\emph{extending} per-client codeword pools for high-uncertainty clients: it deliberately
maintains \emph{divergent}, client-specific codebooks, whereas we keep a \emph{single}
globally-coherent codebook whose index~$j$ has one meaning everywhere---a requirement
forced by our downstream prior, which scores anomalies from the code indices
themselves. In generative and communication settings the codebook is instead
federated by \emph{weight averaging or exchange}: SynthFed~\cite{jiang2026synthfed}
federates a VQ-VAE (with a GPT prior) for privacy-preserving rare-class image
synthesis; FC-DeepSC~\cite{park2022fcdeepsc} and SVQ-VAE~\cite{jin2025svqvae} average
or exchange VQ codebooks to build a shared vocabulary for multi-user deep source
coding and vehicular semantic communication; and UVeQFed~\cite{shlezinger2021uveqfed}
uses vector quantization to \emph{compress model updates}, not to learn a discrete
detector. Two things separate all of these from our setting. First, \emph{objective}:
a reconstruction, transmission, or classification loss tolerates codebook
weight-averaging, whereas our detector consumes code \emph{indices}, so the
cross-client index misalignment that weight-averaging induces
(\S\ref{sec:fed-results}) is hypothesized to collapse the prior's
anomaly signal. Second, \emph{coverage}: none addresses the two remaining
detector-specific pieces we need---aggregating the discrete \emph{prior} over tokens,
and calibrating detection thresholds locally versus globally. Our Stage-1 merge from
per-code sufficient statistics (Proposition~1) is exactly the globally-coherent
alternative to codebook weight-averaging that a federated discrete \emph{detector}
requires.

\paragraph{Federated time-series anomaly detection.}
Federated anomaly detection for time series is a comparatively small but active line:
a dedicated benchmark~\cite{liu2022fedtadbench}, federated VAE/reconstruction
detectors~\cite{zhang2021fedanomaly,zhu2022deepfedad}, and industrial-control
cyberattack detection (TemporalFED~\cite{perales2023temporalfed}). None of these
federates a \emph{discrete/VQ} detector---all use continuous-latent autoencoders or
their variants---so the codebook-federation problem we target does not arise for them.
Relative to this line, we differ in two ways and state our status plainly. First,
\emph{modality and object}: we federate a \emph{discrete/VQ} detector on \emph{real
univariate} KPIs (WSD/AnoTransfer), where the code indices are consumed by a downstream
masked prior---the codebook-federation problem these continuous-latent detectors do not
face. Second, \emph{honesty of scope}: we do not claim a federated accuracy gain. On the
real benchmark a short-budget pilot has federation \emph{below} local-only, but every arm
is under-trained, so we pre-register a converged confirmatory experiment
(\S\ref{sec:fed-convergence}, Table~\ref{tab:fed-main-wsd}) rather than draw a method
conclusion. What we establish now is a real-data mechanism and heterogeneity analysis
(\S\ref{sec:fed-results}); a privacy-leakage frontier and contamination/non-IID robustness
are future work.

% ============================================================================
%  BODY of the -U (UNIVARIATE) federated paper.
%  Scope: federating TimeVQVAE-AD across UNIVARIATE clients (one KPI per client).
%  Primary data: wsd_fed (real, AnoTransfer/WSD, 31 univariate KPI clients).
%  Controlled data: toy_fed_uni (14 univariate synthetic clients, jitter heterogeneity).
%  This is a DIFFERENT paper from the multivariate -M draft (preserved in
%  paper_M_backup_20260710/). Do not reintroduce C=8 multivariate synthetic claims here.
%  Honesty rule: the confirmatory convergence experiment is NOT yet run; its rows are
%  pre-registered and marked \pending. Do not fabricate converged numbers.
% ============================================================================

\section{Method}
\label{sec:fed-method}

\subsection{Problem Setup}
\label{sec:fed-setup}
Each client $s\in\{1,\dots,S\}$ holds a single \emph{univariate} series (one KPI,
$C{=}1$) and never shares raw signals. Locally, TimeVQVAE-AD learns a discrete
time--frequency representation with a VQ-VAE (Stage~1) and fits a MaskGIT masked prior
over the resulting code indices (Stage~2); detection combines a reconstruction term
$s_{\mathrm{local}}$ and a prior masked-NLL term $s_{\mathrm{prior}}$. Because the series
are univariate, the codebook $\mathcal{E}=\{e_j\}_{j=1}^{K}\subset\mathbb{R}^{d}$ is
quantized over a single channel; the per-channel positional embedding degenerates to one
vector. Clients are heterogeneous: they differ in base period, waveform morphology, noise,
and anomaly mix, and this heterogeneity survives per-entity standardization. The question
is what, if anything, is worth federating across such clients, and how to aggregate it.

\subsection{Threat Model and Privacy as a Measured Property}
\label{sec:fed-threat}
Raw signals and latent tokens never leave a client. Federation communicates only
aggregate model state: for Stage~1, per-code assignment counts and centroid-sums
(\S\ref{sec:fed-stage1}); for Stage~2, the transformer body of the prior
(\S\ref{sec:fed-stage2}). We report the per-round upload surface (floats per client per
round) as a measured quantity, not a claim, and defer a formal privacy--leakage frontier
to future work.

\subsection{Federated Stage One: the Codebook as Federated Prototypes (Contribution A)}
\label{sec:fed-stage1}
The codebook is not a gradient-trained parameter. Its assignment
$q_{b,h,w}=\arg\min_j\|z_{b,h,w}-e_j\|_2^2$ is non-differentiable, and each step
overwrites the table from EMA sufficient statistics, $e_j\leftarrow M_j/\tilde{N}_j$, where
$\tilde{N}_j$ is the quantizer's Laplace-smoothed count (Eq.~\ref{eq:fed-round-stats}); we write
$M_j/N_j$ informally throughout, but the implementation smooths the denominator and
Remark~\ref{rem:smoothing} bounds the resulting perturbation. Two
consequences shape how it should be federated. First, the codebook is the sole component
whose shape is independent of the client (a fixed table of $K$ vectors in
$\mathbb{R}^{d}$), hence the natural cross-client object to share. Second, because the
codebook is a set of EMA centroids rather than a loss-trained tensor, averaging its
\emph{weights} (\textsc{FedAvg}) is the wrong primitive; the correct one aggregates the
\emph{sufficient statistics} that define those centroids.

\paragraph{Aggregation rule.} Let $\mathcal{E}^{(t)}$ be the global codebook broadcast at
round $t$. Each client runs $E$ local epochs with the codebook \emph{frozen} to
$\mathcal{E}^{(t)}$ (the encoder/decoder are updated by the straight-through estimator and
the commitment loss) and accumulates \emph{decay-free} round statistics against the shared
assignments:
\begin{equation}
  n^{s}_{j}=\big|\mathcal{Z}^{s}_{j}\big|,\quad
  m^{s}_{j}=\!\!\sum_{z\in\mathcal{Z}^{s}_{j}}\!\!z,\quad
  \mathcal{Z}^{s}_{j}=\{z^{(s)}_{b,h,w}:q^{(s)}_{b,h,w}=j\},
\end{equation}
with $q^{(s)}$ computed against the frozen $\mathcal{E}^{(t)}$. The server then performs
the round's pooled M-step directly from the raw statistics, with the same Laplace-style
smoothing used by the quantizer:
\begin{align}
  N_j &= \textstyle\sum_{s} n^{s}_{j}, &
  M_j &= \textstyle\sum_{s} m^{s}_{j},
  \label{eq:fed-round-stats}\\[2pt]
  \tilde N_j &= \tfrac{N_j+\epsilon}{\sum_i N_i+K\epsilon}\textstyle\sum_i N_i, &
  e^{(t+1)}_j &\leftarrow M_j/\tilde N_j .
  \label{eq:fed-centroid}
\end{align}
Codebook initialization and dead-code revival happen \emph{only} on the server, acting on
globally under-used codes ($N_j<\tau$), which preserves the index correspondence that
\eqref{eq:fed-round-stats} requires.

\begin{proposition}[Aggregation correctness of the federated M-step]
\label{prop:exact}
Fix the broadcast codebook $\mathcal{E}^{(t)}$ and let $\mathcal{A}=\{(z,j)\}$ be the set
of (latent, assigned-index) pairs produced across all clients during the round against
$\mathcal{E}^{(t)}$. Then
$\displaystyle \bar e_j=\frac{\sum_{s} m^{s}_{j}}{\sum_{s} n^{s}_{j}}$
equals the centroid of $\{z:(z,j)\in\mathcal{A}\}$: the per-code mean a single machine
holding $\mathcal{A}$ would compute. The sufficient-statistic merge is an exact, lossless
aggregation of the round's assignments.
\end{proposition}
\begin{proof}[Proof sketch]
Given the shared frozen $\mathcal{E}^{(t)}$, the nearest-neighbour map is identical on
every client, so $\{\mathcal{Z}^{s}_{j}\}_s$ partition
$\mathcal{Z}_j=\bigcup_s\mathcal{Z}^{s}_{j}$; hence $\sum_s m^{s}_{j}=\sum_{z\in\mathcal{Z}_j}z$
and $\sum_s n^{s}_{j}=|\mathcal{Z}_j|$, whose ratio is the pooled centroid. $\square$
\end{proof}

\begin{remark}[Smoothed denominator]
\label{rem:smoothing}
The implementation applies the quantizer's Laplace smoothing to the denominator, so the
merged centroid is $M_j/\tilde{N}_j$ rather than $M_j/N_j$. The perturbation relative to the
pooled centroid is $|N_j/\tilde{N}_j-1|\approx\epsilon/N_j$: measured $10^{-5}$ at $N_j{=}1$,
$10^{-6}$ at $N_j{=}10$, below $10^{-8}$ at $N_j{=}1000$. For well-supported codes the merge
therefore agrees with the pooled M-step to fp32 rounding ($\approx 1.5\times10^{-7}$,
one ulp), which is the tolerance the unit test asserts.
\end{remark}

\begin{remark}[This is the distributed Lloyd M-step, applied to a VQ codebook]
Proposition~\ref{prop:exact} is not a new theorem, and its precedent is older and closer than
we previously stated. It is the standard \emph{distributed} $k$-means M-step of
Dhillon \& Modha~\cite{dhillon1999parallel}, whose Fig.~2 uses the same per-cluster
$(n_j, m_j)$ accumulators, the same pair of sum-allreduces and the same smoothed denominator;
equivalently it is the VQ-VAE EMA update~\cite[App.~A.1, Eqs.~6--8]{van_den_oord2017neural}
with the accumulation boundary moved from the minibatch to the client round. Its
secure-aggregation form is~\cite{diaa2025fastlloyd}; the one-shot
variant k-FED~\cite{dennis2021kfed} is a different algorithm and is \emph{not} the
precedent for this message. We claim no mathematical novelty for the operator.
%
% What we do NOT claim any more (2026-07-27 audit): that this is "the correct primitive for
% merging a VQ-VAE codebook". Dhillon & Modha's descent guarantee assumes every processor
% applies the SAME distance map in a fixed R^d, whereas our E-step is a per-client LEARNED
% encoder that drifts every round -- so no single objective is being descended, and the
% exactness, while real, is vacuous here. The measurable statement that replaces it is the
% negative one: the merged codebook is exact to fp32 while the induced client token languages
% have (measured) disjoint support. Keep that framing.
Our contribution is (i) the observation that weight-space aggregation is the wrong primitive
for this component---prior federated VQ work either averages or exchanges the codeword table
as ordinary model weights~\cite{park2022fcdeepsc,jin2025svqvae}, or sidesteps the
index-alignment problem entirely by maintaining divergent per-client
codebooks~\cite{zhang2024uefl}---and (ii) an empirical characterization of \emph{when} it
matters on real univariate KPIs.
\end{remark}

\begin{remark}[Scope of Proposition~\ref{prop:exact}]
The proposition concerns the \emph{aggregation operator}, not the equivalence of the
federated and centralized \emph{processes}. Encoders are local and update within the
round, so $\mathcal{A}$ differs from what a single centralized encoder would produce;
cross-round semantic stability of the codes is an \emph{empirical} property, encouraged by
the ordinary VQ commitment loss against the frozen broadcast codebook and measured by
token-assignment agreement (\S\ref{sec:fed-metrics}). Summing already-decayed per-client
EMA buffers would \emph{not} reproduce the pooled M-step; this is why clients upload raw
round statistics.
\end{remark}

\begin{remark}[Why not \textsc{FedAvg}/\textsc{FedProx}/\textsc{SCAFFOLD} on the codebook]
Averaging per-client centroids $\tfrac{1}{S}\sum_s e^{s}_j$ equals the pooled centroid only
when $n^{s}_j$ is constant across clients; otherwise low-support (noisy) clients dominate
it. \textsc{FedProx} and \textsc{SCAFFOLD} do not apply to the codebook at all---no loss
gradient reaches $\mathcal{E}$---and are relevant only to the gradient-trained components.
\end{remark}

\subsection{Federated Stage Two: Partially-Personalized Prior (Contribution B)}
\label{sec:fed-stage2}
The MaskGIT prior is federated by \textsc{FedAvg}-ing its transformer body while keeping a
small set of parameters \emph{local}: the per-position output bias (and, in the general
multivariate code, the channel positional embedding, which for univariate clients is a
single vector). Local heads let each client keep an idiosyncratic marginal over code usage
while sharing the sequence model. We ablate this against a fully-global prior
(\S\ref{sec:fed-arms}). We propose no new personalization mechanism; the contribution is
the specific split for this discrete-token prior and its effect on worst-client detection.

\subsection{Is This Federated Learning?}
\label{sec:fed-isfl}
Contribution~A uploads per-code counts and centroid-sums rather than a weight tensor,
which invites the objection that Stage~1 is not federated learning at all. We answer it on
four grounds, and then concede the part of it that stands.

\paragraph{The setting is defined by data locality, not by the message format.}
Federated learning is defined as ``a machine learning \emph{setting} where many clients
\dots\ collaboratively train a model under the orchestration of a central server \dots\
while keeping the training data decentralized''~\cite{kairouz2021advances};
\textsc{FedAvg}~\cite{mcmahan2017fedavg} is one instantiation of that setting, not its
definition. The constraint that defines it---raw signals never leave a client, a server
aggregates---holds throughout our protocol (\S\ref{sec:fed-threat}). Aggregating
sufficient statistics rather than weights is ordinary practice inside it, and federated
$k$-means is the direct precedent: its per-round protocols sum per-cluster counts and
point-sums across clients---structurally the message we upload---and recent work layers
secure aggregation and differential privacy on exactly that
pair~\cite{diaa2025fastlloyd}. The one-shot variant, k-FED, instead has each device solve a
local $k$-means problem and communicate its local cluster \emph{means}, which the server
re-clusters~\cite{dennis2021kfed}. Both designs must solve the same problem---making index
$j$ denote the same thing on every client---and they solve it differently: k-FED by
re-clustering the centroids after the fact, ours by anchoring them to a frozen common
broadcast beforehand, which is what makes the per-round merge exact
(Prop.~\ref{prop:exact}). We do not run k-FED as an arm because its regime does not hold
here: its guarantee assumes each device holds data from only $k'\le\sqrt{k}$ of the $k$
target clusters, whereas a codebook is sized so that a client uses a large fraction of it
(\S\ref{sec:fed-metrics} reports the measured codebook perplexity), and because it is
one-shot over fixed points whereas ours are encoder outputs that move during training. What
the centroid-message design costs in-regime is instead characterized analytically in
\S\ref{sec:fed-stage1}: averaging per-client centroids recovers the pooled centroid only
when the supports $n^{s}_{j}$ are equal across clients. More broadly, federated
analytics computes counts and rates over decentralized data under secure aggregation with
no gradient-based training anywhere in the loop~\cite{ramage2020fedanalytics,bonawitz2022privacy}.

Closest to our message---because it aggregates statistics of \emph{latent} vectors rather
than of raw points---is federated prototype learning. \textsc{FedProto} defines a client's
per-class prototype as the mean of its embedding vectors, $m^{s}_{j}/n^{s}_{j}$ in our
notation, and the server combines them with count weights $n^{s}_{j}/\sum_s n^{s}_{j}$,
which is the pooled centroid~\cite[Eqs.~3,~6]{tan2022fedproto}. Its stated rationale is
the one we inherit: under heterogeneity ``the optimal model parameters for each client are
not the same[, which] means that gradient-based communication cannot sufficiently provide
useful information to each client,'' whereas a shared label space lets clients ``share the
same embedding space'' so that information can be exchanged by aggregating per-class
statistics. Our setting substitutes the codeword index for the class label, and the frozen
broadcast codebook does the work FedProto's shared label space does: it is what makes index
$j$ denote the same thing on every client. We upload the raw pair $(n^{s}_{j}, m^{s}_{j})$
rather than the ratio because the server needs both to form the pooled centroid, and
because a client may have $n^{s}_{j}=0$ for codes it never used.

\paragraph{For a codebook, the sufficient statistics \emph{are} the model update.}
The objection presumes that a model update is a weight delta. For a VQ table it is not.
Already in the original VQ-VAE the codeword has a closed-form optimum equal to the mean of
the encoder outputs assigned to it---an update its authors describe as ``typically used in
algorithms such as $K$-means''---and the practical minibatch rule is an exponential moving
average over precisely the pair $(N_i, m_i)$ of accumulated counts and accumulated sums,
with $e_i = m_i/N_i$~\cite[App.~A.1]{van_den_oord2017neural}. Uploading $(n^{s}_{j}, m^{s}_{j})$
therefore transmits the codebook update in its natural coordinates, and
$e_j = M_j/\tilde N_j$ (Eq.~\ref{eq:fed-centroid}) is a deterministic decoding of that
message rather than a different object. A weight-space message is not a richer alternative
but a lossy re-encoding of the same content: it discards the supports $n^{s}_{j}$, and
averaging the resulting per-client centroids recovers the pooled centroid only when those
supports are equal across clients (\S\ref{sec:fed-stage1}). Proposition~\ref{prop:exact} is
the consequence: the statistic-space merge is exact, and the weight-space one approximates it.

\paragraph{The message format is not a privacy deficit.}
The uploaded statistics are \emph{linear} in the client's latents, so the operation the
server requires is a plain sum and is compatible with secure aggregation by
construction~\cite{bonawitz2017secure}; under it the server observes only
$(\sum_s n^{s}_{j}, \sum_s m^{s}_{j})$ and no per-client message. The weight-space
alternative that the objection prefers is, conversely, the one message type demonstrably
invertible to raw inputs~\cite{zhu2019dlg,geiping2020inverting}. We claim no privacy
guarantee here---\S\ref{sec:fed-threat} treats leakage as a measured surface and defers a
formal frontier to future work---only that ``it shares no weights'' is not by itself a
deficiency.

\paragraph{The objection is answered empirically, and by Contribution~B.}
The protocol is not weight-free: Stage~2 federates the MaskGIT transformer body by
\textsc{FedAvg} (\S\ref{sec:fed-stage2}), so the objection can apply at most to Stage~1.
And the weight-space federation of Stage~1 that it asks for is not omitted---it is
arm~(4)~\emph{--A} (\S\ref{sec:fed-arms}), run under the same budget, seeds and selection
rule. Whether the statistic-space primitive earns its place is therefore an empirical
question we answer rather than assume; at $n{=}31$ this particular contrast is underpowered
and is reported as such (\S\ref{sec:fed-protocol}).

\paragraph{What we concede.}
With encoder and decoder kept local, Stage~1 federates a single $K\times d$ table---a small
fraction of the parameters---which is closer to \emph{agreeing on a shared vocabulary} than
to collaborative representation learning. We do not dispute that reading; we report the
per-round upload surface as a measured quantity so that it can be judged
(\S\ref{sec:fed-threat}). But it is the finding rather than an evasion of it: the question
posed in \S\ref{sec:fed-setup} is what, if anything, is worth federating across
heterogeneous univariate clients, and a small answer is still an answer. We test the
shared-vocabulary reading directly rather than leaving it rhetorical, via cross-client
token agreement on a shared probe and worst-client detection
(\S\ref{sec:fed-metrics}): a merely better-estimated quantizer would improve the former
without redistributing the latter.

\subsection{Local Calibration and Detection}
\label{sec:fed-detect}
After training, each client fixes an operating threshold as the $q{=}0.99$ quantile of its
own \emph{train} scores (test is never used for calibration). On the univariate KPI data
the deployed score is the prior term alone ($s_{\mathrm{local}}$ weight $0$), so detection
performance is a function of $s_{\mathrm{prior}}$; we exploit this in the mechanism
analysis (\S\ref{sec:fed-results}).

\section{Experimental Design}
\label{sec:fed-experiments}

\subsection{Research Questions}
\begin{description}
  \item[\textbf{RQ1 (Viability).}] At a converged budget, does the federated recipe reach
  local-only detection \emph{without sharing raw data}, per client and worst-case, and does
  it approach the centralized (data-pooling) reference?
  \item[\textbf{RQ2 (Primitive).}] Does sufficient-statistic aggregation
  \eqref{eq:fed-round-stats}--\eqref{eq:fed-centroid} beat \textsc{FedAvg}-on-codebook-weights,
  and through which score term?
  \item[\textbf{RQ3 (What must be personalized).}] Does a partially-personalized prior
  preserve worst-client detection better than a fully-global prior?
  \item[\textbf{RQ4 (When federation hurts).}] Which client property predicts the
  federated--local gap on real KPIs---cohort heterogeneity, or the client's own
  detectability?
\end{description}

\subsection{Datasets and Client Construction}
% CITATION BUG (was \cite{dennis2021kfed} = k-FED clustering, which does NOT supply WSD).
% \bibitem{anotransfer} MUST BE ADDED -- no AnoTransfer entry exists in the bibliography.
% Provenance of the frozen subset: dataset-federated/MODEL.md + dataset-federated/scripts/build_frozen.py.
\textbf{Real benchmark (primary): \textsc{wsd\_fed}.} We use WSD, the AnoTransfer KPI
corpus~\cite{anotransfer}: real univariate web-service KPIs at 1-minute sampling with a
24-hour period and $\sim$0.6\% point anomalies. Our frozen subset has $31$ clients grouped
into $4$ cohorts by train-only KMeans on an average-day deviation descriptor
($k{=}4$ is the most seed-stable choice; silhouette $0.53$; membership is train-only, no
test leakage). Per-client training windows range $2734$--$16541$ ($W{=}256$, stride $1$).
% Rebuilt 2026-07-16 from 14 -> 47 clients (2 clients/cluster was too few to federate).
% Per-client series lengths unchanged; entity ids were REUSED (uni_04 was M2_valve, is now
% M1_rotary), so pre-rebuild checkpoints load silently and give wrong results. Every toy
% number below predates the rebuild and must be re-measured before it is cited.
\textbf{Controlled benchmark (secondary): \textsc{toy\_fed\_uni}.} $47$ univariate
synthetic clients in $6$ machine-type cohorts (sizes $6/7/9/11/8/6$), with an explicit within-cohort
heterogeneity axis (\texttt{jitter}$\in\{$mild, medium, strong$\}$) and injected anomalies,
used to probe the mechanism under controlled non-IID structure.

\subsection{Methods Compared}
\label{sec:fed-arms}
Two reference arms bracket the range: (1)~\textsc{Local-only} (one model per client, no
FL) and (2)~\textsc{Centralized} (one model on pooled windows; a data-sharing reference
trained for the same number of \emph{epochs} as local, hence with more optimizer steps in
proportion to cluster size --- the contrast is confounded by compute and we treat it as an
upper bracket, not a controlled contrast (\S\ref{sec:fed-protocol})). Our recipe and its ablations isolate each contribution:
(3)~\textsc{Ours\,(A$+$B)}, the sufficient-statistic codebook merge with the
partially-personalized prior; (4)~\emph{--A}, a weight-\textsc{FedAvg} codebook with the
same partial personalization; (5)~\emph{--B}, the sufficient-statistic codebook with a
fully-global prior. Personalized-FL methods (FedBN, Ditto, pFedMe) and external federated
time-series AD baselines are future work. A commitment-weight ``anchor'' knob was tested
and \emph{dropped}: in federated mode it reduces algebraically to a rescaling of the
commitment weight (the anchor term $\lambda\,\ell_{\mathrm{commit}}$ adds to
$w_{\mathrm{c}}\,\ell_{\mathrm{commit}}$), so it is not an independent regularizer and we do
not report it as a mechanism.

\subsection{Metrics and Primary Endpoint}
\label{sec:fed-metrics}
\textbf{Primary (pre-specified).} VUS-PR~\cite{boniol2022vus}, range-aware and endorsed as
robust for TSAD; the tolerance buffer is derived from the data ($14$ for \textsc{wsd\_fed}).
\textbf{Secondary.} AUPRC, AUROC, PATE, affiliation-F1. \textbf{Legacy only.} thresholded
F1 / point-adjusted F1, reported for comparability but \emph{not} used to draw conclusions,
since point adjustment and thresholded F1 are known to inflate scores~\cite{kim2022towards}.
Because tolerance buffers differ across datasets, VUS-PR is \emph{not} compared across
\textsc{wsd\_fed} and \textsc{toy\_fed\_uni}; each dataset is analyzed on its own.
\textbf{Mandatory calibration baseline.} a centered moving-average residual (a one-line
detector) is pre-registered to be reported alongside the deep model on \textsc{wsd\_fed}, so
that absolute numbers are interpretable; this baseline is \pending\ (\S\ref{sec:fed-results}).
\textbf{Mechanism.} we decompose detection into $s_{\mathrm{local}}$ (reconstruction) and
$s_{\mathrm{prior}}$ (masked-NLL) and report each separately, plus codebook health
(perplexity, dead-code fraction) and cross-client token agreement on a shared probe.

\subsection{Evaluation Protocol (Statistical Plan)}
\label{sec:fed-protocol}
All arms share fixed hyper-parameters and a compute schedule; the validation split is a
disjoint temporal tail of train, used only for training diagnostics and \emph{symmetric}
model selection across all arms (no arm gets a validation-restore the others lack). The
% federated_eval.py:205 (per-client loader) vs :216 (pooled loader), both with s1_epochs.
% No step-matching code exists on this path; the centralized arm is epoch-matched only.
centralized arm is \emph{epoch}-matched to local, not step-matched: it trains for the same
number of epochs over the pooled cluster, hence for $3.4$--$11.4\times$ the gradient steps of
the median client (\S\ref{sec:fed-results}). The centralized--local contrast therefore
confounds data pooling with compute and is reported as a bracket, not a controlled contrast.
\textbf{Unit of analysis.} we report per-entity
paired contrasts with the \emph{entity} as the unit ($n{=}31$ on \textsc{wsd\_fed}), and we
disclose that entities within a cohort share one federated model, so entity-level
$p$-values are mildly anticonservative; we therefore also report cohort-level
sign-consistency. At $n{=}31$ the minimum detectable effect (paired Wilcoxon, $80\%$ power)
is $\approx 0.087$ VUS-PR, so effects below that---plausibly including the
suff-stat-vs-\textsc{FedAvg} contrast---are underpowered on this benchmark and are reported
as such. We use $\geq 4$ seeds and report wall-clock and GPU-hours.

\subsection{Convergence Requirement (pre-registered)}
\label{sec:fed-convergence}
A pilot at a short budget ($3$ rounds / $3$ epochs) leaves every arm far from convergence
(federated Stage-2 validation still falling $\sim$40\% at the last round; the smallest
client sees only $\sim$130 Stage-1 gradient steps against the $\sim$1500 of a converged
run). We therefore \emph{pre-register} the confirmatory budget: $\geq 15$ rounds /
$\geq 30$ epochs (with per-client epochs scaled so the smallest client reaches
$\gtrsim 1500$ Stage-1 steps), and symmetric validation selection. (The centralized arm
remains epoch-matched, not step-matched --- no step-matching protocol has been implemented.)
All RQ1--RQ3 conclusions are drawn only at this converged budget; short-budget
results are reported as a pilot, not as evidence about the method.

\section{Results}
\label{sec:fed-results}

% Converged sweep is MEASURED on disk (artifacts/fed_eval/converged/wsd_fed.csv: 744 rows =
% 6 arms x 31 entities x 4 seeds, meta 35 ep / 18 rounds x 2 local ep / seeds 0-3, meeting the
% pre-registration at sec:fed-convergence) but is NOT yet transcribed into this table pending a
% final audit. "Not yet reported here", never "not yet measured".
\paragraph{Status.} The confirmatory converged sweep (\S\ref{sec:fed-convergence}) is measured
but not yet reported here; the rows below marked \pending\ are pre-registered and awaiting a
final audit. We report here (i) the pilot viability numbers with their undertraining caveat,
and (ii) the mechanism and heterogeneity analyses that are already available from the pilot
artifacts on real data. The calibration baseline is \pending\ (\S\ref{sec:fed-metrics}).

% WITHDRAWN 2026-07-16. This paragraph asserted the deep model loses to a 10-min MA on the
% strength of local 0.306 vs MA 0.330. BOTH numbers were defective: 0.330 has no artifact (see
% the table comment) and 0.306 is the UNDER-TRAINED pilot; converged local on disk
% (artifacts/fed_eval/converged/wsd_fed.csv, 35 ep / 4 seeds) is markedly higher. The random
% ~0.01 / "50-100x lift" and "beats seasonal-1440/z-score/long-window MA" claims are likewise
% unbacked (no such baseline artifact exists). The sign of this comparison is UNKNOWN until the
% MA baseline is run and matched against CONVERGED local.
\paragraph{Calibration --- pending.} A mandatory calibration baseline (a centred moving-average
residual, \S\ref{sec:fed-metrics}) is pre-registered and \pending. We therefore make no claim
about the deep detector's absolute advantage in either direction; all federation claims are
\emph{relative} (arm vs arm under one protocol) by construction, never absolute SOTA.

\begin{table}[t]
\centering\small
\caption{\textsc{wsd\_fed} detection (VUS-PR, per-entity mean over $n{=}31$; entity is the
unit). Pilot rows are at the short $3/3$ budget and are \emph{under-trained}; the
\pending\ rows are the pre-registered converged sweep (\S\ref{sec:fed-convergence}). No
converged federated numbers are claimed yet.}
\label{tab:fed-main-wsd}
\begin{tabular}{lcc}
\toprule
Method & Pilot ($3/3$, under-trained) & Converged (pre-registered) \\
\midrule
% 0.330 REMOVED 2026-07-16: no code and no artifact on disk produces it. There is no standalone
% moving-average detector in this repo (no scripts/*, no _arm in any records_wsd_fed.jsonl);
% detect.py _moving_average_paper is the (a+MA)/2 impulse term of the deep score, not a baseline.
% Its only source is prose in documentation/DEAD_AND_OPEN_LINES.md (section O5, listed as NOT YET RUN).
% TO FILL: implement a centred 10-min MA residual detector, score all 31 wsd_fed entities with the
% same VUS-PR (buffer 14) used for every other row, and write it to artifacts/fed_eval/ as a real arm.
Moving-average (10\,min)            & \pending & \pending \\
Local-only (no FL)                  & 0.306 & \pending \\
% NOT step-matched, at any budget: federated_eval.py:205 trains each local model s1_epochs on its
% own loader; :216 trains centralized s1_epochs on the POOLED loader (3.4-11.4x the median client's
% steps, tracking cluster size). "Epoch-matched" is the only accurate word.
Centralized (pooled data, epoch-matched)\textsuperscript{$\dagger$} & 0.335 & \pending \\
\textbf{Ours (A$+$B)}               & 0.115 & \pending \\
\emph{--A}: FedAvg-weights codebook & \pending & \pending \\
\emph{--B}: global prior            & \pending & \pending \\
\bottomrule
\end{tabular}
\\[2pt]
% run_fed_converged.sh passes one --s1-epochs 35 to BOTH arms; federated_eval.py has no step
% budget. A genuinely step-matched centralized arm has never been run.
{\footnotesize \textsuperscript{$\dagger$}Centralized is \emph{not} step-matched
($3.4$--$11.4\times$ more gradient steps than the median client); the converged sweep on disk
uses the same epoch-matched protocol and does \emph{not} fix this.}
\end{table}

\paragraph{RQ1 (pilot, not confirmatory) --- under-trained federation loses to local.}
At the short budget the federated recipe is the worst of the three arms in all four cohorts
(per-entity $\Delta_{\text{fed}-\text{local}}={-}0.19$ VUS-PR, $n{=}31$, sign-consistent in
$27/31$ entities). We do \emph{not} interpret this as a property of the method: every arm is
far from convergence (\S\ref{sec:fed-convergence}). The converged sweep is required to make
RQ1 interpretable.
% REMOVED 2026-07-16: "the suff-stat-vs-FedAvg contrast is exactly zero at this short budget"
% has no artifact on disk backing an EXACT zero, and every toy_fed_uni number predates the
% 2026-07-16 14->47 rebuild. Re-measure on the 47-client build before restating.

\paragraph{RQ2/RQ4 mechanism --- federation degrades the prior, not reconstruction; and on
real data the driver is detectability, not heterogeneity.}
The deployed \textsc{wsd\_fed} score is the prior term alone ($s_{\mathrm{local}}$ weight
$0$), so any federated collapse is, by construction, a \emph{prior} collapse; the Stage-1
codebook stays healthy under federation (final-round perplexity $20$--$34/64$, dead codes
driven to $0\%$ by server revival). This is verifiable per-entity from the persisted
checkpoints without retraining (autopsy, CPU). On the \emph{controlled} synthetic
% STALE + STATISTICALLY FRAGILE (re-measure before citing): all these toy numbers predate the
% 2026-07-16 14->47 rebuild. The p=1e-4 is on entity x seed pseudo-replicates (n=42), not the
% n=14 entities the paper elsewhere declares the unit; and the strong-jitter level (-0.34) rests
% on 2 entities. Re-run on the 47-client build with the entity as the unit.
\textsc{toy\_fed\_uni}, the federated--local gap is monotone in within-cohort heterogeneity
(\texttt{jitter}: mean $\Delta$ ${-}0.01/{-}0.09/{-}0.34$; Spearman ${-}0.56$; on the retired
$14$-client build, pending re-measurement). On \emph{real}
\textsc{wsd\_fed}, however, that heterogeneity signal is \emph{null}
(descriptor-distance vs gap $\rho{=}{+}0.10$, $p{=}0.61$); the only strong predictor is the
client's own local score ($\rho{=}{-}0.77$, $p{=}4\mathrm{e}{-}7$): federation acts as
regression-to-the-mean, pulling strong clients down and weak clients up. We report the
synthetic-vs-real discrepancy as a finding, not paper over it.

\paragraph{RQ2/RQ3 primitive and personalization --- pre-registered.}
The confirmatory contrasts (suff-stat vs weight-\textsc{FedAvg}; partial vs global prior)
require the converged sweep and are \pending\ on \textsc{wsd\_fed}. On controlled synthetic
data the mechanism signature we expect is a prior-side collapse under weight-\textsc{FedAvg}
(token-assignment agreement falling toward the uniform-random floor while reconstruction
stays intact); confirming or refuting this on real KPIs is the load-bearing experiment.

\paragraph{Scope and honest negatives.}
% (i) previously asserted the MA loss on the strength of 0.306 vs 0.330; 0.330 has no artifact
% and 0.306 is the pilot. Withdrawn 2026-07-16 -- see the table comment.
(i) The pre-registered moving-average calibration baseline (\S\ref{sec:fed-metrics}) is
\pending, so the detector's absolute advantage is neither claimed nor refuted here. (ii) Under-trained federation loses to local; whether
this survives at convergence is the open question the pre-registered sweep resolves. (iii)
The suff-stat primitive's advantage, seen on multivariate synthetic data in our own earlier
experiments (not an external prior result), is \emph{untested} on real univariate data and did
not separate from \textsc{FedAvg} at short budgets on the (now-retired 14-client) univariate
synthetic build---so it is a hypothesis here, not a result. (iv) Cohort heterogeneity does
not predict the federated gap on real data. A privacy--leakage frontier and
contamination/non-IID robustness are out of scope.


\section{Conclusion and Limitations}
\label{sec:fed-conclusion}

We studied how to federate TimeVQVAE-AD across univariate KPI clients that keep raw
signals local, while making the discrete time--frequency representation and its masked
prior trainable across clients. The recipe has two parts: \textbf{(A)} a global VQ
codebook merged from per-code sufficient statistics $(n^{s}_{j}, m^{s}_{j})$, whose
server-side merge reproduces the pooled $k$-means M-step (federated $k$-means applied to a
VQ codebook, not a new theorem; Prop.~\ref{prop:exact}, \S\ref{sec:fed-stage1}); and
\textbf{(B)} a partially-personalized MaskGIT prior whose transformer body is
\textsc{FedAvg}'d while small per-client output heads stay local (\S\ref{sec:fed-stage2}).

Rather than overclaim, we report an honest, in-progress picture on the \emph{real}
$31$-client univariate KPI benchmark (\S\ref{sec:fed-results}).
% withdrawn: "The deep detector does not beat a moving-average baseline" -- unmeasured, see
% federated_method.tex:253. Also removed "exactly zero" (unbacked) and the stale toy claim.
All federation comparisons are relative (arm vs arm under one protocol) by construction. A
short-budget pilot has the federated recipe losing to local-only in every cohort, but every
arm is far from convergence---so the pilot is not evidence about the method; we
pre-register the converged confirmatory sweep. What we \emph{can} establish now is
mechanistic and about \emph{when} federation hurts: the deployed KPI score is the prior
term alone, so federated degradation is a \emph{prior} collapse with a healthy Stage-1
codebook; and the ``heterogeneity drives the federated gap'' law that holds on synthetic
data is \emph{null} on real data, where the only predictor is a client's own
detectability (federation as regression-to-the-mean).

\paragraph{Limitations.} First, \emph{the confirmatory experiment is pending}: RQ1
(parity) and RQ2 (the sufficient-statistic primitive vs weight-\textsc{FedAvg}) are drawn
only at a converged budget, which the pilot does not reach; we therefore make no method
claim from the pilot. Second, \emph{absolute quality is unmeasured}: the pre-registered
% was "the deep detector ties a 10-min moving-average" -- "ties" and "does not beat" are
% different claims from the same absent measurement (0.306 vs 0.330 is a loss, not a tie).
moving-average calibration baseline (\S\ref{sec:fed-metrics}) has not been run, so we claim
only relative, protocol-matched contrasts. Third, \emph{statistical power is bounded}: with $31$ entities the paired
minimum detectable effect is $\approx 0.087$ VUS-PR, so a small suff-stat advantage would
be undetectable here, and entities within a cohort share a model, making entity-level
$p$-values mildly anticonservative. Fourth, we run \emph{no privacy or contamination
experiments}: the upload surface is characterized qualitatively (\S\ref{sec:fed-threat}),
not measured for leakage.

\paragraph{Future work.} The immediate next step is the pre-registered converged sweep
(\S\ref{sec:fed-convergence}) that resolves RQ1--RQ3 on real KPIs; alongside it, a
\emph{damped} exact merge ($\alpha\,e^{\text{local}}{+}(1{-}\alpha)\,e^{\text{merged}}$) as
a fix for the short-horizon reversal of the primitive's advantage---the exact M-step being
the right estimator but a poor optimizer step before encoders align. Beyond that: an
empirical \emph{privacy--leakage frontier} on the uploaded statistics, paired with secure
aggregation~\cite{bonawitz2017secure} and DP; a heterogeneity-controlled univariate
benchmark with more high-heterogeneity clients; and testing whether the
codebook-versus-prior mechanism is a reusable primitive for federating other
discrete-representation detectors.

\begin{thebibliography}{99}
\bibitem{mcmahan2017fedavg}
McMahan, B., Moore, E., Ramage, D., Hampson, S., Ag\"uera y Arcas, B.:
Communication-Efficient Learning of Deep Networks from Decentralized Data.
In: AISTATS, pp.\ 1273--1282 (2017)

\bibitem{li2020fedprox}
Li, T., Sahu, A.K., Zaheer, M., Sanjabi, M., Talwalkar, A., Smith, V.:
Federated Optimization in Heterogeneous Networks. In: MLSys (2020)

\bibitem{tan2022fedproto}
Tan, Y., Long, G., Liu, L., Zhou, T., Lu, Q., Jiang, J., Zhang, C.:
FedProto: Federated Prototype Learning across Heterogeneous Clients.
In: AAAI, vol.\ 36(8), pp.\ 8432--8440 (2022)

\bibitem{dennis2021kfed}
Dennis, D.K., Li, T., Smith, V.:
Heterogeneity for the Win: One-Shot Federated Clustering.
In: ICML, PMLR 139, pp.\ 2611--2620 (2021)

\bibitem{li2021fedbn}
Li, X., Jiang, M., Zhang, X., Kamp, M., Dou, Q.:
FedBN: Federated Learning on Non-IID Features via Local Batch Normalization.
In: ICLR (2021)

\bibitem{li2021ditto}
Li, T., Hu, S., Beirami, A., Smith, V.:
Ditto: Fair and Robust Federated Learning Through Personalization.
In: ICML, pp.\ 6357--6368 (2021)

\bibitem{dinh2020pfedme}
T.\ Dinh, C., Tran, N.H., Nguyen, T.D.:
Personalized Federated Learning with Moreau Envelopes. In: NeurIPS (2020)

\bibitem{zhu2019dlg}
Zhu, L., Liu, Z., Han, S.: Deep Leakage from Gradients. In: NeurIPS (2019)

\bibitem{bonawitz2017secure}
Bonawitz, K., et al.: Practical Secure Aggregation for Privacy-Preserving
Machine Learning. In: ACM CCS, pp.\ 1175--1191 (2017)

\bibitem{so2023securing}
So, J., Ali, R.E., G\"uler, B., Jiao, J., Avestimehr, A.S.:
Securing Secure Aggregation: Mitigating Multi-Round Privacy Leakage in
Federated Learning. In: AAAI (2023)

\bibitem{kairouz2021advances}
Kairouz, P., McMahan, H.B., et al.:
Advances and Open Problems in Federated Learning.
Found.\ Trends Mach.\ Learn.\ 14(1--2), 1--210 (2021)

\bibitem{boniol2022vus}
Paparrizos, J., Boniol, P., Palpanas, T., Tsay, R.S., Elmore, A., Franklin, M.J.:
Volume Under the Surface: A New Accuracy Evaluation Measure for Time-Series
Anomaly Detection. In: VLDB (2022)

\bibitem{huet2022local}
Huet, A., Navarro, J.M., Rossi, D.:
Local Evaluation of Time Series Anomaly Detection Algorithms. In: ACM SIGKDD (2022)

\bibitem{kim2022towards}
Kim, S., Choi, K., Choi, H.-S., Lee, B., Yoon, S.:
Towards a Rigorous Evaluation of Time-series Anomaly Detection. In: AAAI (2022)

\bibitem{fourure2021how}
Fourure, D., Javaid, M.U., Posocco, N., Tihon, S.:
Anomaly Detection: How to Artificially Increase Your F1-Score with a Biased
Evaluation Protocol. In: ECML PKDD (2021)

\bibitem{perales2023temporalfed}
S\'anchez S\'anchez, P.M., Huertas Celdr\'an, A., Bovet, G., Mart\'inez P\'erez, G.:
TemporalFED: Detecting Cyberattacks in Industrial Time-Series Data Using
Decentralized Federated Learning. arXiv:2308.03554 (2023)

\bibitem{kapoor2024reforms}
Kapoor, S., Cantrell, E., Peng, K., et al.:
REFORMS: Consensus-based Recommendations for Machine-learning-based Science.
Science Advances 10(18) (2024)

% NOTE: the EMA codebook update over $(N_i,m_i)$ cited in \S\ref{sec:fed-isfl} is
% Eqs.\ (6)--(8) of Appendix~A.1 of the arXiv version ONLY; the NeurIPS proceedings text
% mentions moving averages in passing and says they are "not used for the experiments in
% this work". The arXiv pointer below is therefore load-bearing, not decorative.
\bibitem{van_den_oord2017neural}
van den Oord, A., Vinyals, O., Kavukcuoglu, K.:
Neural Discrete Representation Learning. In: NeurIPS, vol.\ 30, pp.\ 6306--6315 (2017).
EMA codebook update over $(N_i,m_i)$: Appendix~A.1 of arXiv:1711.00937

\bibitem{lee2023timevqvae}
Lee, D., Malacarne, S., Aune, E.:
Vector Quantized Time Series Generation with a Bidirectional Prior Model.
In: AISTATS, PMLR 206, pp.\ 7665--7693 (2023)

\bibitem{lee2024timevqvaead}
Lee, D., Malacarne, S., Aune, E.:
Explainable Time Series Anomaly Detection using Masked Latent Generative Modeling.
Pattern Recognition 156, 110826 (2024)

\bibitem{chang2022maskgit}
Chang, H., Zhang, H., Jiang, L., Liu, C., Freeman, W.T.:
MaskGIT: Masked Generative Image Transformer. In: CVPR, pp.\ 11315--11325 (2022)

\bibitem{arivazhagan2019fedper}
Arivazhagan, M.G., Aggarwal, V., Singh, A.K., Choudhary, S.:
Federated Learning with Personalization Layers. arXiv:1912.00818 (2019)

\bibitem{collins2021fedrep}
Collins, L., Hassani, H., Mokhtari, A., Shakkottai, S.:
Exploiting Shared Representations for Personalized Federated Learning.
In: ICML, pp.\ 2089--2099 (2021)

\bibitem{pillutla2022partial}
Pillutla, K., Malik, K., Mohamed, A., Rabbat, M., Sanjabi, M., Xiao, L.:
Federated Learning with Partial Model Personalization.
In: ICML, pp.\ 17716--17758 (2022)

\bibitem{zhang2024uefl}
Zhang, T., Cao, Y., Liu, D.:
Uncertainty-Based Extensible Codebook for Discrete Federated Learning in
Heterogeneous Data Silos. In: ICML (2025). arXiv:2402.18888

\bibitem{jiang2026synthfed}
Jiang, X., Qian, J., Gu, A., Ma, X., Jin, K., Zhang, X., Song, Z.:
SynthFed: Privacy-Preserving Long-Tail Ophthalmic Diagnosis via VQ-VAE and
GPT-Augmented Federated Learning. Biomedical Signal Processing and Control 113 (2026)

\bibitem{park2022fcdeepsc}
Park, C.-H., Choi, J., Park, J., Kim, S.-L.:
Federated Codebook for Multi-User Deep Source Coding.
In: ICTC, pp.\ 994--996 (2022)

\bibitem{jin2025svqvae}
Jin, Z., Li, Y., Song, T., Jia, W.-K., Song, X.:
SVQ-VAE: Federated-Learning-Based Semantic-Aware Communication for Vehicular
Networks. IEEE Internet of Things Journal 12(23), 51289--51304 (2025)

\bibitem{shlezinger2021uveqfed}
Shlezinger, N., Chen, M., Eldar, Y.C., Poor, H.V., Cui, S.:
UVeQFed: Universal Vector Quantization for Federated Learning.
IEEE Trans.\ Signal Processing 69, 500--514 (2021)

\bibitem{lu2023hvqtrans}
Lu, R., Wu, Y., Tian, L., Wang, D., Chen, B., Liu, X., Hu, R.:
Hierarchical Vector Quantized Transformer for Multi-class Unsupervised Anomaly
Detection. In: NeurIPS (2023)

\bibitem{trappolini2025qaed}
Trappolini, G., Purificato, A., Siciliano, F., D'Addona, L., Spagnolo, A.M.,
Dato, D., Silvestri, F.:
Quantized Auto Encoder-Based Anomaly Detection for Multivariate Time Series Data
in 5G Networks. IEEE Access (2025)

\bibitem{liu2022fedtadbench}
Liu, F., Zeng, C., Zhang, L., Zhou, Y., Mu, Q., Zhang, Y., Zhang, L., Zhu, C.:
FedTADBench: Federated Time-Series Anomaly Detection Benchmark. In: IEEE HPCC (2022)

\bibitem{zhang2021fedanomaly}
Zhang, K., Jiang, Y., Seversky, L., Xu, C., Liu, D., Song, H.:
Federated Variational Learning for Anomaly Detection in Multivariate Time Series.
In: IEEE IPCCC (2021)

\bibitem{zhu2022deepfedad}
Zhu, W., Song, D., Chen, Y., Cheng, W., Zong, B., Mizoguchi, T., Lumezanu, C.,
Chen, H., Luo, J.:
Deep Federated Anomaly Detection for Multivariate Time Series Data. arXiv:2205.04041 (2022)

% --- added 2026-07-27 for \S sec:fed-isfl; all fields verified against primary sources ---

\bibitem{diaa2025fastlloyd}
Diaa, A., Humphries, T., Kerschbaum, F.:
FastLloyd: Federated, Accurate, Secure, and Tunable $k$-Means Clustering with
Differential Privacy. In: USENIX Security Symposium, pp.\ 2733--2752 (2025)

\bibitem{ramage2020fedanalytics}
Ramage, D., Mazzocchi, S.:
Federated Analytics: Collaborative Data Science without Data Collection.
Google Research Blog, 27 May 2020

\bibitem{geiping2020inverting}
Geiping, J., Bauermeister, H., Dr\"oge, H., Moeller, M.:
Inverting Gradients --- How Easy Is It to Break Privacy in Federated Learning?
In: NeurIPS, vol.\ 33, pp.\ 16937--16947 (2020)

\bibitem{bonawitz2022privacy}
Bonawitz, K., Kairouz, P., McMahan, B., Ramage, D.:
Federated Learning and Privacy. Communications of the ACM \textbf{65}(4), 90--97 (2022)

\bibitem{anotransfer}
Zhang, S., Zhong, Z., Li, D., Fan, Q., Sun, Y., Zhu, M., Zhang, Y., Pei, D.,
Sun, J., Liu, Y., Yang, H., Zou, Y.:
Efficient KPI Anomaly Detection Through Transfer Learning for Large-Scale Web
Services. IEEE Journal on Selected Areas in Communications \textbf{40}(8),
2440--2455 (2022)
\end{thebibliography}

\end{document}
